\documentclass[12pt,a4paper]{article}

\usepackage{graphicx}
\usepackage{enumitem}
\usepackage{hyperref}
\usepackage[english,iso]{isodate}

\newcommand{\authorname}{Joe J Collins}
\newcommand{\documenttitle}{The Curation of Empirical Data}
\makeatletter
\newcommand{\subtitle}[1]{\g@addto@macro\@title{\\[0.5em]\Large #1}}
\makeatother

\hypersetup{
  pdftitle={\documenttitle},
  pdfauthor={\authorname}
}

\title{\documenttitle}
\subtitle{A Seven-Layer Architecture for Maintaining an Unbroken Chain of Provenance in Field Research}
\author{\authorname}
\date{\today}

\begin{document}

\maketitle

\begin{abstract}
  Write a concise summary of the article here.
\end{abstract}

\tableofcontents

\clearpage
%%%%%%%%%%%%%%%%%%%%%%%%
\section{Seven Steps}

Typically,
the collection, cleaning and curation of data is not formally taught
to students of the sciences.
There is enough for students to learn in their own disciplines
without having to learn the intricacies of data management.
As a result, many students develop their own ad hoc methods for
collecting, cleaning and curating data.
Most often
the student will determine the shape of the data frame required for their analysis.
Then take this data frame
(in the form of a spreadsheet)
to the field and attempt to fill it with data.
This approach is fraught with problems of data integrity and provenance.
At best it is inefficient and error-prone.
However,
for a single student working alone on a small project,
it may be sufficient.
However, this approach is
not scalable or sustainable
once the student joins a research team
and is no longer the sole collector, cleaner and consumer of the data.

Fortunately, the process of collecting, cleaning and curating data
has a long history in the sciences and engineering.
The process is often poorly documented
and often obscured by the technical details of the software used to implement it.
For software vendors there is a perverse incentive to obscure the process
to make it appear that their software is doing something magical.
In reality, the process is often simple and straightforward.
The purpose of this article is to describe a simple and straightforward process
for collecting, cleaning and curating data.

Each step in the process is a separate layer in a seven-layer architecture
for maintaining an unbroken chain of provenance in field research data.
Focusing on them separately allows explicit control over the process
is reproducible and transparent.
Separating the layers allows for the process to be scaled and sustained
across multiple researchers and multiple projects.
And allows researchers to focus effectively on each step in the process
without being distracted by the other steps.

%%%%%%%%%%%%%%%%%%%%%%%%
\section{Collect}

\subsection{Measure}

\subsection{Record}

\subsection{Share}

%%%%%%%%%%%%%%%%%%%%%%%%
\section{Enhance}

\subsection{Refine}

\subsection{Regulate}

\subsection{Retrieve}

\subsection{Wrap}

%%%%%%%%%%%%%%%%%%%%%%%%
\section{Investigate}




\section*{References}

\addcontentsline{toc}{section}{References}

\end{document}